% !PRINTED 654
\weberpage{654}

\begin{center}
\S.\,181.\\
\gesperrt{Vorbereitung zum Beweis.}
\end{center}

Nach den Ausführungen des vorigen Paragraphen ist zum Beweis des grossen Theorems I nur noch der Nachweis erforderlich, dass die in $\Omega_m$ enthaltene Zahl
\[
\omega = (r, x_0)^m
\]
die $m^{\text{te}}$ Potenz einer Kreistheilungszahl ist.

1. \gesperrt{Die wesentliche Eigenschaft der zu $\omega$ conjugirten Zahlen $\omega_n$, auf die sich der Beweis stützt, ist die, dass}
\weq{1}{\omega_n^{-1} \omega_1^n = \alpha^m}
\gesperrt{die $m^{\text{te}}$ Potenz einer Zahl in $\Omega_m$ ist.}

Hierin kann $n$ jede durch $q$ nicht theilbare Zahl sein. Bilden wir von der rechten und linken Seite von (1) die Normen im Körper $\Omega_m$, so ergiebt sich, wenn wir die Norm der Zahl $\alpha$, die eine rationale Zahl ist, mit $a$ bezeichnen, da jedes $\omega_n$ dieselbe Norm wie $\omega$ hat,
\[
N(\omega)^{n-1} = a^m.
\]
Wenn nun $m$ ungerade ist, so können wir hierin $n = 2$ annehmen, und erhalten als \gesperrt{Folgerung aus 1.}:

2. \gesperrt{Die Norm von $\omega$ ist die $m^{\text{te}}$ Potenz einer rationalen Zahl}
\weq{2}{N(\omega) = a^m.}

% [sic] The German source prints the anomalous punctuation ``Potenz von 2.,''; retained verbatim.
Ist $m$ eine Potenz von 2., so ist diese Folgerung nicht mehr zulässig. Es ergiebt sich dann aus 1. nur, dass $N(\omega)^2$ die $m^{\text{te}}$ Potenz einer rationalen Zahl ist. Gleichwohl müssen die Zahlen $\omega$, wie aus (10) des vorigen Paragraphen zu sehen ist, auch in diesem Falle noch der Bedingung 2. genügen, die dann aber nicht mehr von selbst schon in 1. enthalten ist, sondern als neue Eigenschaft hinzukommt.

\gesperrt{Was wir zu beweisen haben, ist, dass aus 1. und 2. folgt, dass auch $\omega$ selbst die $m^{\text{te}}$ Potenz einer Kreistheilungszahl, wenn auch in einem höheren Körper, ist.}

Dazu ist aber erforderlich, die Zerlegung der Zahlen $\omega$ in ihre Primfactoren im Körper $\Omega_m$ zu untersuchen.

Die Primzahl $q$ ist, wie im \S.\,169 nachgewiesen, mit der $\varphi(m)^{\text{ten}}$ Potenz einer in $\Omega_m$ existirenden Primzahl $\sigma = 1 - r$
%CONT>>
